\documentclass[11pt]{amsart}
\usepackage{amsmath,amssymb,amsthm}
\usepackage[margin=1in]{geometry}
\usepackage{microtype}
\usepackage{hyperref}

\theoremstyle{plain}
\newtheorem{thm}{Theorem}[section]
\newtheorem{lem}[thm]{Lemma}
\newtheorem{prop}[thm]{Proposition}
\newtheorem{cor}[thm]{Corollary}
\newtheorem{conj}[thm]{Conjecture}
\theoremstyle{definition}
\newtheorem{defn}[thm]{Definition}
\newtheorem{rem}[thm]{Remark}

\newcommand{\Ch}{\operatorname{Ch}}
\newcommand{\cube}{\square}
\newcommand{\atom}[1]{\langle #1 \rangle}

\title[The structure of the Dedekind endgame]
{The structure of the Dedekind endgame: one approximation statement and
 two auxiliaries}

\author{Ryota Shibusawa}
\address{Daiichi Institute of Technology, Japan}
\email{r-shibusawa@daiichi-koudai.ac.jp}

\subjclass[2020]{55U35, 18N40; 18G55, 06D30}
\keywords{cubical set, Dedekind cube, chain approximation, chain stratum,
  subatom order, type/test model structure}

\begin{document}

\begin{abstract}
  The comparison of the type and test model structures on the Dedekind
  cubical site was reduced, in the companion manuscript, to a single
  cubical approximation statement together with two combinatorial
  statements about the subatom order.  We record a structural
  clarification that unifies the three: both combinatorial statements
  feed the \emph{same} chain-approximation master statement.  The first,
  well-foundedness of the subatom order, reduces to a tight bounded-%
  generation conjecture (companion note).  The second,
  intrinsic-chain-anchored asphericity, is shown here to be
  \emph{equivalent} to chain approximation on the anchored subobjects of
  atoms: the intrinsic chain stratum of an anchored subobject $W$ with
  $\Ch(A)\subseteq W\subseteq A$ satisfies $\Ch(W)=\Ch(A)$, so $W$ is
  type-contractible iff its own chain approximation holds.  Thus the
  Dedekind endgame has a single homotopical core---chain
  approximation---and two auxiliaries: a finiteness (securing the
  well-founded induction) and the anchored instance of the core itself.
  The anchored instance is test-verified in every computed case.
\end{abstract}

\maketitle

\section{The endgame after reduction}

In the companion manuscript the Dedekind type/test comparison is proved
equivalent, through Cisinski regularity, a cosimplicial resolution and a
singular Quillen adjunction, to a single \emph{chain-approximation}
statement: for every fibrant $X$ the inclusion of the intrinsic chain
stratum $\Ch(X)\hookrightarrow X$ is a type equivalence.  The singular
complex computes the \emph{test} homotopy type of every presheaf, so
$\Ch(X)\hookrightarrow X$ is always a test equivalence; only the
type-level upgrade remains.  An adversarially audited assembly then
localizes that upgrade to two combinatorial statements:
\begin{itemize}
\item[(a)] the subatom order of $\cube^n$ is well-founded;
\item[(b)] every subobject $W$ of a cyclic atom $A$ (or its fold and
  stabilizer quotients) with $\Ch(A)\subseteq W\subseteq A$ is
  type-contractible (\emph{intrinsic-chain-anchored asphericity}).
\end{itemize}
We record how (a) and (b) both reduce to the chain-approximation core.

\section{(a) reduces to bounded generation}

\begin{thm}[companion note]
  The measure $|{\operatorname{im}}|$ is non-increasing along the subatom
  order, and if for each finite lattice $L$ there are finitely many atoms
  with image $L$ then the order is well-founded.  This finiteness follows
  from the tight conjecture that an atom is generated at arity at most the
  height of its image; the conjecture is verified, and tight, in every
  computed case.
\end{thm}

So (a) secures the well-founded induction that drives the assembly; it is
a \emph{finiteness}, orthogonal to the homotopy content.

\section{(b) is the anchored instance of chain approximation}

Write $\Ch(A)$ for the intrinsic chain stratum of an atom $A$: the
sorted cells of $A$ and their chain instances (equivalently the image of
$\Ch_n$ under the sort action), which is a subpresheaf.  A subobject $W$
is \emph{anchored} if $\Ch(A)\subseteq W\subseteq A$.

\begin{lem}\label{lem:chw}
  For an anchored subobject $W$, $\Ch(W)=\Ch(A)$.
\end{lem}

\begin{proof}
  Sortedness is a property of a cell itself, independent of the ambient
  object; hence the sorted cells of $W$ are exactly the sorted cells of
  $A$ that lie in $W$.  Since $\Ch(A)\subseteq W$, every sorted cell of
  $A$ and every chain instance of one already lies in $W$, so the sorted
  cells of $W$ are all of those of $A$ and their chain instances again
  lie in $W$.  Therefore $\Ch(W)=\Ch(A)$.  (Machine-checked on the
  median atom and others: \texttt{dedekind\_o32\_anchor.py}.)
\end{proof}

\begin{prop}\label{prop:bequiv}
  Let $A$ be a cyclic atom, so $\Ch(A)\hookrightarrow A$ is a type
  equivalence and $\Ch(A)$ is type-contractible.  For an anchored $W$,
  \[
    W \text{ is type-contractible}
    \iff
    \Ch(W)\hookrightarrow W \text{ is a type equivalence.}
  \]
  Hence (b) holds iff chain approximation holds on every anchored
  subobject of every cyclic atom---a sub-case of the master statement.
\end{prop}

\begin{proof}
  By Lemma~\ref{lem:chw}, $\Ch(W)=\Ch(A)$, which is type-contractible.
  If $\Ch(W)\hookrightarrow W$ is a type equivalence then $W\simeq\Ch(W)$
  is contractible.  Conversely, if $W$ is contractible then both $W$ and
  $\Ch(W)=\Ch(A)$ are contractible, so the inclusion between them is a
  type equivalence.
\end{proof}

\begin{cor}
  For an anchored $W$, the inclusion $\Ch(W)\hookrightarrow W$ is always
  a \emph{test} equivalence (the singular complex computes the test type,
  and $\Ch(W)=\Ch(A)$ is test-contractible), so $W$ is test-contractible.
  The content of \textup{(b)} is exactly the type-level upgrade---the
  same type-versus-test gap as the master statement, restricted to
  anchored subobjects.
\end{cor}

Thus (b) is not an independent obstruction: it is chain approximation
itself, evaluated on the anchored subobjects of atoms.

\section{The unified picture}

\[
  \boxed{\;
  \begin{array}{c}
  \textbf{Master: chain approximation on all fibrant }X
  \iff \text{Dedekind comparison}\\[4pt]
  \text{(a) well-foundedness}\ \Longleftarrow\ \text{bounded generation (finiteness)}\\[2pt]
  \text{(b)}\ \Longleftrightarrow\ \text{chain approximation on anchored }W\subseteq\text{ atoms}
  \end{array}
  \;}
\]

The Dedekind endgame therefore has a single homotopical core, chain
approximation, and two auxiliaries of different character: (a) is a
finiteness that makes the assembling induction well-founded, and (b) is
the master statement's own restriction to the anchored subobjects of
atoms.  The cone mechanism that contracts constant-free atoms does not
extend to this restriction---it fails already on the median atom---so a
non-cone, triangulation-level contraction is what remains.

\section{Verification}

For the median atom (where the cone fails) and the constant-free
five-chain, and across a sweep of sixteen atoms with $|L|\in\{3,4,5\}$,
median- and non-median-fibre alike, every sampled anchored subobject $W$
is test-contractible (mod-$2$ Betti $1,0,0$); $\Ch(A)$ is a proper
subobject in the majority of cases, so the tests are not vacuous.  The
scripts are \texttt{dedekind\_o32\_conechain.py} (cone (non-)preservation
of the intrinsic stratum) and \texttt{dedekind\_o32\_anchor.py}
(anchored-$W$ Betti sweep), in the artifact of the companion manuscript.

\begin{thebibliography}{9}
\bibitem{PaperW}
R.~Shibusawa. Which cubical sites present spaces? Manuscript, 2026.
\bibitem{NoteWF}
R.~Shibusawa. The well-foundedness of the cubical subatom order reduces
to bounded generation arity. Note, 2026.
\end{thebibliography}

\end{document}
